\section{Research Design and Methodology}

To achieve these aims, our project is organized into two primary research strands running in parallel. This methodology combines direct spatial profiling of human tissue with live, functional assays in vitro.

\subsection{Key Project Strands}

The project is structured around two distinct yet complementary research strands designed to run in parallel.

\noindent
\begin{minipage}[t]{0.48\textwidth}
    \begin{tcolorbox}[
        enhanced,
        title=\sffamily\bfseries Strand A: Spatial Profiling,
        colback=bgcolor,
        colframe=primarycolor,
        coltitle=white,
        arc=4pt,
        boxrule=1.5pt,
        height=4.0cm
    ]
        \small
        This strand focuses on the mapping of transcriptomic changes directly in human post-mortem tissues. By using high-resolution multiplexed error-robust fluorescence in situ hybridization (MERFISH), we will construct a high-fidelity reference map of human synaptic RNA localization.
    \end{tcolorbox}
\end{minipage}%
\hfill
\begin{minipage}[t]{0.48\textwidth}
    \begin{tcolorbox}[
        enhanced,
        title=\sffamily\bfseries Strand B: Functional Validation,
        colback=bgcolor,
        colframe=secondarycolor,
        coltitle=white,
        arc=4pt,
        boxrule=1.5pt,
        height=4.0cm
    ]
        \small
        This strand leverages advanced microfluidic models of synaptic circuits. Using these platforms, we will perform live-imaging and localized translation assays to validate the functional roles of target transcriptomic variants discovered in Strand A.
    \end{tcolorbox}
\end{minipage}

\subsection{Strand A: High-Resolution Spatial Transcriptomics}
We will analyze human post-mortem cortical tissues obtained from the Meridian Brain Bank. We will deploy MERFISH, targeting a panel of 500 candidate synaptic transcripts selected based on preliminary single-nucleus RNA-sequencing data. Using sub-diffraction imaging and advanced optical decoding, we will identify the spatial coordinates of these RNAs relative to pre- and post-synaptic markers (e.g., Synaptophysin and PSD-95). This allows us to quantify changes in RNA abundance and localization at the single-synapse level.

\subsection{Strand B: Microfluidic Modeling of the Synaptic Environment}
In collaboration with Nexa Biotech, we will utilize state-of-the-art multi-compartment microfluidic devices that isolate axons and synapses from cell bodies. Induced pluripotent stem cells (iPSCs) from patients carrying familial Alzheimer's mutations will be cultured in these chambers. A custom microfluidic platform designed by Nexa Biotech allows for 95\% transfection efficiency within isolated compartments, enabling us to selectively manipulate local transcript levels and measure real-time changes in synaptic transmission and morphological remodeling.

\begin{figure}[htbp]
    \centering
    \begin{tikzpicture}[
        node distance=1.8cm,
        box/.style={
            draw=primarycolor,
            fill=bgcolor,
            thick,
            rounded corners=4pt,
            text width=3.8cm,
            align=center,
            font=\sffamily\small,
            minimum height=1.2cm
        },
        arrow/.style={
            -stealth,
            thick,
            color=secondarycolor
        }
    ]
        % Nodes
        \node[box] (obj1) {\textbf{Objective 1: Spatial Mapping}\\MERFISH profiling of synaptic microdomains.};
        \node[box, right=of obj1, xshift=1.2cm] (obj2) {\textbf{Objective 2: Functional Assays}\\Microfluidic platforms to test local translation.};
        \node[box, below=of obj1, yshift=-0.5cm] (obj3) {\textbf{Objective 3: Drug Screening}\\In vitro screening of Synapse Inc. candidate therapeutics.};
        \node[box, right=of obj3, xshift=1.2cm] (obj4) {\textbf{Objective 4: Translation}\\In vivo validation with Apex Labs translational pipeline.};
        
        % Arrows
        \draw[arrow] (obj1) -- (obj2);
        \draw[arrow] (obj2) -- (obj4);
        \draw[arrow] (obj1) -- (obj3);
        \draw[arrow] (obj3) -- (obj4);
        
        % Background
        \begin{scope}[on background layer]
            \draw[draw=secondarycolor!20, fill=secondarycolor!2, dashed, rounded corners=8pt]
                ($(obj1.north west) + (-0.5, 0.5)$) rectangle ($(obj4.south east) + (0.5, -0.5)$);
            \node[anchor=south west, font=\sffamily\bfseries\color{secondarycolor}] at ($(obj1.north west) + (-0.5, 0.45)$) {Research Design \& Translation Pathway};
        \end{scope}
    \end{tikzpicture}
    \caption{Overview of the proposed research framework spanning spatial profiling, functional validation, drug screening, and translation.}
    \label{fig:workflow}
\end{figure}

\subsection{Project Milestones \& Key Deliverables}
To ensure rigorous monitoring of the project's progress, we have established clear milestones across the five-year funding period (Table~\ref{tab:milestones}).

\begin{table}[htbp]
\centering
\sffamily\small
\renewcommand{\arraystretch}{1.2}
\begin{tabularx}{\textwidth}{l l X}
\toprule
\rowcolor{secondarycolor}
\textbf{\color{white}Milestone} & \textbf{\color{white}Target Month} & \textbf{\color{white}Description} \\
\midrule
M1.1 & Month 12 & Complete MERFISH spatial mapping of 50 post-mortem human brain samples. \\
\rowcolor{bgcolor}
M2.1 & Month 24 & Establish microfluidic synapse models with 95\% transfection efficiency. \\
M3.1 & Month 36 & Validate candidate synaptic protective compounds in microfluidic models. \\
\rowcolor{bgcolor}
M4.1 & Month 48 & Translate top 3 lead compounds into preclinical in vivo models with Apex Labs. \\
M5.1 & Month 60 & Publish open-access spatial transcriptomic database of the human synapse. \\
\bottomrule
\end{tabularx}
\caption{Key milestones and timeline targets for the proposed grant.}
\label{tab:milestones}
\end{table}
